\documentclass{article}
\usepackage{iclr2026_conference,times}

\input{math_commands.tex}

\usepackage{hyperref}
\usepackage{url}
\usepackage{amsmath}
\usepackage{amssymb}
\usepackage{booktabs}
\usepackage{graphicx}
\usepackage{xcolor}
\usepackage{amsthm}
\newtheorem{proposition}{Proposition}

% ---- drafting aid: remove before submission -------------------------------
\newcommand{\todo}[1]{\textcolor{red}{\textbf{[TODO: #1]}}}
\newcommand{\pending}[1]{\textcolor{blue}{\textit{[pending: #1]}}}
% ---------------------------------------------------------------------------

\title{{\sc PACLock}: Gauge-Invariant Phase--Amplitude\\
Tokenization for EEG}

% ICLR 2026 提交必须匿名：\iclrfinalcopy 保持注释状态。
% 模板原文：Non-anonymous submissions will be rejected without review.
% 录用后取消下面两处注释即可切换到 camera-ready。
\author{Zhizhe Zhang \\
Department of Computer Science \\
New York University \\
\texttt{zz5070@nyu.edu}}

%\iclrfinalcopy % 仅 camera-ready 时取消注释，提交时必须保持注释

\begin{document}

\maketitle

\begin{abstract}
Neurophysiology gives EEG a strong structural prior: slow oscillatory
\emph{phase} gates fast oscillatory \emph{amplitude}, a phenomenon known as
phase--amplitude coupling (PAC). Yet architectures that inject this prior as a
learnable module --- an additive coupling bias, a gate, a PAC-weighted attention
---consistently fail to separate from plain attention. We show why: in every such
design the prior sits \emph{beside} a free path that the optimiser can take
instead, and it does. We therefore make the prior \emph{constitutive} rather than
modulatory. In our operator the token for a fast band is not a representation of
that band which PAC then reweights; it \emph{is} the product of that band's
amplitude with the preferred-phase-aligned sum of all slower bands' phases. Above
the slowest band no raw path exists, so the prior cannot be routed around. The
construction has a further property no prior PAC model has: it is provably
invariant to the arbitrary choice of phase reference, because the measured
preferred phase cancels exactly the gauge rotation it induces on the learned
phase features. We further give an evaluation design whose parameter-matched
controls are constructed to isolate the measured preferred phase from the three
explanations that would otherwise confound it: the multiplicative form itself,
the mere exposure of analytic phase features, and weighting by coupling strength.
\pending{empirical results; downstream task suite and baseline set are not yet
fixed and are deliberately omitted from this draft}
\end{abstract}

\section{Introduction}

Foundation models for EEG \citep{yang2023biot,jiang2024labram,wang2025cbramod}
tokenize the signal along time and space: a recording becomes a grid of
(electrode, time-patch) tokens, and any frequency structure the model uses must
be recovered by the encoder. This is a deliberate design choice --- it makes the
tokenizer domain-agnostic --- but it discards a prior that neuroscience regards
as one of the most robust organising principles of cortical dynamics.

That prior is \emph{phase--amplitude coupling} (PAC): the amplitude envelope of a
fast oscillation is modulated by the phase of a slower one
\citep{canolty2006high,tort2010measuring}. PAC is directional (slow gates fast),
it is measurable per electrode and per short time window, and it carries two
distinct quantities --- a coupling \emph{strength} and a \emph{preferred phase},
the phase of the slow rhythm at which the fast envelope peaks.

A natural idea is to build this into the architecture. Several works do
\citep{lemkhenter2020phaseswap,li2023pacnet,li2023complexpac,lee2026sleeppacnet},
and so did we, repeatedly. Our central negative finding is that a large class of
these designs fails for a shared and predictable reason.

\paragraph{Priors with an escape route get optimised away.}
We ran four architectural injections of PAC on identical backbones and data
(Section~\ref{sec:whyfail}). An additive $\alpha \cdot \mathrm{PAC}$ bias drove
its scalar $\alpha \to 0$. A sigmoid gate, initialised at its no-op value so it
could ``only help,'' never left that value over the whole of training --- an
entire result table that turned out to measure plain attention. Two hyperparameter-free
PAC-modulated attention operators, one soft and one a hard top-$k$ topology,
matched but did not beat plain attention, and did not beat their own
label-destroying shuffle controls. In each case the module sat next to a path
that could carry the same information without it, and training preferred that
path.

\paragraph{Our design rule.}
If a prior can be bypassed, it will be. So we do not add PAC to a token; we
\emph{define} the token with it. For a band $j$ above the slowest, the token is
the elementwise product of that band's amplitude feature with a
preferred-phase-aligned, coupling-weighted sum of the phase features of all
slower bands. There is no raw waveform token beside it, no scale, no gate, no
auxiliary branch. Every sample-dependent feature that reaches the encoder for
$j>0$ contains both factors, inseparably.

\paragraph{Gauge invariance.}
The phase origin of a band is arbitrary: change the reference electrode or the
filter implementation and every phase shifts by a constant. A representation
built on phase should not depend on that choice. Our operator satisfies this
\emph{exactly and by construction} (Section~\ref{sec:gauge}): a shift $\delta_i$
of band $i$'s phase reference rotates the learned phase feature by
$e^{i\delta_i}$ and rotates the measured coupling vector by the same
$e^{i\delta_i}$, and the operator applies the conjugate of the latter to the
former, so the two cancel. To our knowledge no prior PAC architecture has, or
claims, this property; a survey of 28 related papers returns no occurrence of
phase-reference invariance or equivariance in this sense.

\paragraph{Contributions.}
\begin{enumerate}
\item A documented account of \emph{why} architecturally injected PAC priors
fail, from four controlled negative results, and the design rule this implies:
the prior must be constitutive, not modulatory (Section~\ref{sec:whyfail}).
\item {\sc PACLock}, a mandatory PAC tokenization operator, with a proof of exact
invariance to the phase reference (Sections~\ref{sec:method}--\ref{sec:gauge}).
\item An evaluation design that isolates the responsible quantity rather than
merely demonstrating an aggregate gain. Three parameter-matched controls remove,
respectively, the measured coefficients, the preferred-phase alignment, and the
forced multiplication, so that each competing explanation is falsified separately
(Section~\ref{sec:controls}), together with a capacity control that separates the
operator's effect from model size (Section~\ref{sec:capacity}).
\end{enumerate}

\section{Related work}

\paragraph{EEG foundation models.}
BIOT \citep{yang2023biot} tokenizes per channel with an FFT and pretrains
contrastively. LaBraM \citep{jiang2024labram} trains a vector-quantized tokenizer
against the Fourier spectrum and then predicts masked codes. CBraMod
\citep{wang2025cbramod} reconstructs masked raw patches with criss-cross
spatial/temporal attention. All three tokenize time patches; none carries an
explicit frequency axis, and none computes an interaction between bands as part
of tokenization.

\paragraph{PAC in deep learning.}
PAC is not new to deep models, and we do not claim it is. Phase-Swap
\citep{lemkhenter2020phaseswap} builds a pretext task from exchanging Fourier
phase between segments. PACNet \citep{li2023pacnet} uses event-related PAC to
\emph{select} subject-specific high-frequency bands ahead of an EEGNet-style
decoder. \citet{li2023complexpac} assemble a complex-valued PAC image from Tort's
modulation index and preferred phase and classify it with a complex CNN.
SleepPACNet \citep{lee2026sleeppacnet} is architecturally closest: it Hilbert
transforms a low band, concatenates its real and imaginary parts with a
high-band signal, and lets a CNN discover the interaction.

Our distinction from these is not the use of PAC but three structural properties:
PAC is computed \emph{differentiably inside the network} from a learnable
filterbank rather than precomputed offline; the interaction is \emph{forced}
rather than made available; and the result is \emph{gauge-invariant}. Notably
SleepPACNet retains a separate raw-EEG branch alongside its PAC branch --- by our
analysis exactly the escape route that lets such priors be ignored --- and its
reported gain is $2.4$ accuracy points with only one sleep stage significant
after correction. We test the concatenative alternative directly inside our own
architecture in Section~\ref{sec:controls}.

\section{Why architecturally injected PAC priors fail}
\label{sec:whyfail}

\todo{condense; currently the fullest version of the argument}

We summarise four attempts, all on the same tri-axial backbone, same data, same
protocol, differing only in how PAC enters.

\begin{enumerate}
\item \textbf{Additive bias.} The frequency-axis mixer adds
$\alpha \cdot C$ to the attention logits, $C$ the measured coupling matrix,
$\alpha$ learned. $\alpha$ decays monotonically toward zero during training.
\item \textbf{Multiplicative gate.} Attention probabilities are multiplied by
$\sigma(w \cdot C)$ and renormalised, with $w$ initialised at $0$ so the module
starts as an exact no-op and can ``only switch gating on if it helps.'' Across
five datasets and six layers the gate mean stayed at $0.5$ to the fifth decimal
--- it never left initialisation. The recorded conclusions from those runs
measured plain attention and were retracted.
\item \textbf{Hyperparameter-free multiplicative modulation.} Post-softmax
attention weights are multiplied by row-mean-normalised coupling and
renormalised. No learnable strength exists, so nothing can decay. It ties plain
attention.
\item \textbf{Hard topology.} Coupling selects the top-$k$ source bands per
target; the rest are masked to $-\infty$. Also ties, and does not separate from a
control in which the coupling matrix is shuffled.
\end{enumerate}

The gate case is the most instructive. A module initialised at a no-op value has
two indistinguishable outcomes: it correctly declined to act, or it never moved.
Without an explicit check that the parameter changed, the design is untestable by
construction. We recommend such a check as standard practice.

The common structure is that in (1)--(4) the encoder retains a complete,
sample-dependent token for each band regardless of the PAC module, so the module
is optional. The design rule that follows is the premise of this paper: to make a
prior load-bearing, remove the alternative.

\section{Method}
\label{sec:method}

\subsection{Tri-axial analytic tokenization}

Let $x \in \mathbb{R}^{B\times C\times T}$ be raw EEG. A learnable sinc
filterbank \citep{ravanelli2018sincnet}, parameterised by band edges rather than
kernel weights, produces $n_b$ band-limited signals per electrode. A
differentiable Hilbert transform yields, per band, a unit-modulus analytic phase
$e^{i\varphi_i(t)}$ and an amplitude envelope $A_j(t)$. Time is divided into
patches of length $L$; all quantities below are computed \emph{within} one
(electrode, patch) pair and we suppress those indices.

\subsection{The measured coupling vector}

For a slow band $i$ and a fast band $j$ we take the debiased mean vector
\citep{canolty2006high},
\begin{equation}
Z_{ij} \;=\; \frac{1}{L}\sum_{t} \bigl(A_j(t) - \bar{A}_j\bigr)\, e^{i\varphi_i(t)} ,
\label{eq:Z}
\end{equation}
retaining it as a complex number: $|Z_{ij}|$ is coupling strength and
$\angle Z_{ij}$ is the preferred phase. Centring $A_j$ removes the bias that a
strictly positive weight would otherwise induce. Crucially we do \emph{not} take
the modulus at this point; discarding $\angle Z_{ij}$ is what left earlier
operators with no phase geometry to steer by.

\subsection{Two feature maps}

A single bias-free convolution $\mathcal{L}$, shared between the real and
imaginary parts, produces a complex phase feature; a separate convolution
$\mathcal{A}$ produces an amplitude feature:
\begin{equation}
p_i \;=\; \mathcal{L}(\cos\varphi_i) + i\,\mathcal{L}(\sin\varphi_i),
\qquad
a_j \;=\; s \odot \mathcal{A}\bigl(\log(1+A_j)\bigr).
\label{eq:feats}
\end{equation}
The absence of a bias in $\mathcal{L}$ and its sharing across the two components
are not incidental --- they are exactly the conditions under which
Section~\ref{sec:gauge} holds. The diagonal $s$ exists to match the parameter
count of the raw-token baseline exactly.

\subsection{The mandatory interaction}

Only edges $i<j$ are admissible: slow phase drives fast amplitude. Coupling
strengths are normalised per target band,
\begin{equation}
\alpha_{ij} \;=\; \frac{|Z_{ij}|}{\sum_{i'<j} |Z_{i'j}|},
\end{equation}
which makes the operator invariant to the overall scale of $Z$ and therefore
free of any temperature or strength hyperparameter. The token is
\begin{equation}
\boxed{\;
h_j \;=\; \underbrace{a_j}_{\text{target amplitude}} \odot
\sum_{i<j} \alpha_{ij}\,
\underbrace{e^{-i\angle Z_{ij}}}_{\text{preferred-phase alignment}}
\underbrace{p_i}_{\text{slower-band phase}} ,
\qquad
h_0 \;=\; a_0 \odot p_0 .
\;}
\label{eq:token}
\end{equation}
The alignment factor is what makes the sum meaningful: different edges peak at
different phases, and adding unaligned complex features would cause them to
cancel. The slowest band has no slower driver and forms the root of the directed
hierarchy. Tokens are mapped to $\mathbb{R}^{2K}$ by interleaving real and
imaginary parts, then given a band embedding and an electrode-coordinate
embedding, and passed to a factorised transformer attending along time, space,
and frequency in turn. The frequency-axis mixer is \emph{plain} attention: PAC
enters the model in exactly one place.

\subsection{No bypass}

For $j>0$ there is no raw token beside $h_j$, no learned coupling scale, no gate,
no PAC branch and no auxiliary loss. The positional embeddings are constants and
cannot carry sample-specific signal around the product. Because $a_j$ and the
aligned phase are combined multiplicatively, no downstream layer can recover
either factor alone.

\section{Exact invariance to the phase reference}
\label{sec:gauge}

\begin{proposition}
Let $\varphi_i \mapsto \varphi_i + \delta_i$ for arbitrary constants $\delta_i$.
If $\mathcal{L}$ is real-linear, bias-free, and shared between the real and
imaginary components as in \eqref{eq:feats}, then $h_j$ is unchanged for all
$j>0$.
\end{proposition}

\begin{proof}
Expanding $\cos(\varphi+\delta)$ and $\sin(\varphi+\delta)$ and using linearity
and the absence of a bias,
\[
\mathcal{L}(\cos(\varphi_i{+}\delta_i)) + i\,\mathcal{L}(\sin(\varphi_i{+}\delta_i))
= e^{i\delta_i}\bigl(\mathcal{L}(\cos\varphi_i) + i\,\mathcal{L}(\sin\varphi_i)\bigr),
\]
so $p_i \mapsto e^{i\delta_i} p_i$. From \eqref{eq:Z}, $Z_{ij} \mapsto e^{i\delta_i} Z_{ij}$,
hence $e^{-i\angle Z_{ij}} \mapsto e^{-i\delta_i} e^{-i\angle Z_{ij}}$, while
$\alpha_{ij}$ is unchanged. The two rotations cancel in every term of
\eqref{eq:token}.
\end{proof}

Numerically, applying independent random shifts to all bands changes
$h_{j>0}$ by at most $8.3\times 10^{-7}$, i.e.\ float32 rounding. The root
token $h_0$ is gauge-\emph{covariant} by definition --- some band must fix the
reference --- and is excluded from the claim.

Two consequences are worth stating. First, this is a property of the operator,
not something training must discover, which distinguishes it from a learned
rotation of feature pairs that have no defined phase origin. Second, it
constrains what a phase-based objective may target: a gauge-\emph{covariant}
quantity such as $\angle Z_{ij}$ itself cannot be predicted from a
gauge-invariant representation, so only invariant targets are well posed.

\section{Evaluation design}
\label{sec:eval}

\pending{The downstream task suite and the baseline set are not yet fixed. This
section therefore states the design and the falsification criteria; the tables
are deliberately left empty rather than populated with results that a change of
task suite would invalidate.}

Our claim is not that a PAC tokenizer improves an aggregate number --- that
claim would be satisfied by any of several mechanisms we consider uninteresting.
The claim is that a \emph{specific measured quantity}, the preferred phase, is
what carries the improvement, and that it does so only when the interaction is
made unavoidable. The design below is built so that each competing explanation
can be rejected on its own.

\subsection{Controlled comparison}

Every comparison holds data, splits, preprocessing, optimiser, schedule and seed
fixed and varies only the tokenizer. \textbf{All arms are constructed to have
identical parameter counts}, so no outcome can be attributed to capacity. The
reference point is a raw-token backbone: the same learnable filterbank, the same
frequency axis, the same encoder, differing only in that each token is a
convolution of one band's filtered waveform rather than the interaction
\eqref{eq:token}.

\subsection{Controls that isolate the mechanism}
\label{sec:controls}

Each control removes exactly one ingredient of \eqref{eq:token} and thereby
targets one alternative explanation.

\begin{table}[t]
\caption{Control arms. Each is parameter-matched to the full operator and
disables a single component, so that the corresponding alternative explanation
can be rejected independently.}
\label{tab:controls}
\begin{center}
\begin{tabular}{lll}
\toprule
Arm & Component removed & Explanation it rejects \\
\midrule
raw baseline & the entire interaction & --- (reference point) \\
\texttt{uniform} & measured $\alpha_{ij}$ \emph{and} $\angle Z_{ij}$ & ``the multiplicative form is what helps'' \\
\texttt{magnitude} & $\angle Z_{ij}$ only, deterministically & ``coupling strength alone suffices'' \\
\texttt{concat} & the forced product only & ``exposing analytic phase suffices'' \\
\texttt{scramble} & the phase-to-edge pairing & ``any complex rotation suffices'' \\
\midrule
\texttt{measured} & nothing & --- \\
\bottomrule
\end{tabular}
\end{center}
\end{table}

Two of these deserve comment. \texttt{uniform} retains the entire multiplicative
topology and replaces only the measured coefficients with a uniform average; if
the multiplicative form were itself responsible, this arm would retain the gain.
\texttt{concat} is the sharper control: it receives \emph{identical ingredients}
--- the same amplitude feature and the same gauge-invariant aligned phase --- and
differs only in combining them through a learned linear map instead of the forced
product. It is the concatenative fusion used by the closest prior work
\citep{lee2026sleeppacnet}, evaluated inside our own architecture so that
exactly one variable changes: whether the interaction is optional. Note that
\texttt{concat} necessarily carries slightly more parameters than the product
arm, which biases the comparison \emph{against} our operator.

\texttt{scramble} preserves every $|Z_{ij}|$ and the entire multiset of preferred
phases but permutes which edge owns which, redrawing the permutation on every
forward pass so that no fixed permutation can be learned and inverted. Because
this injects stochasticity that the deterministic arms lack, we use
\texttt{magnitude}, not \texttt{scramble}, when decomposing the contribution of
the preferred phase.

\begin{table}[t]
\caption{Main comparison. \pending{task suite not fixed}}
\label{tab:main}
\begin{center}
\begin{tabular}{lccc}
\toprule
Task & raw token & PAC interaction & $\Delta$ \\
\midrule
\pending{---} & & & \\
\bottomrule
\end{tabular}
\end{center}
\end{table}

\subsection{Separating the operator from model capacity}
\label{sec:capacity}

A tokenizer that helps a small model may only be compensating for limited
capacity. We therefore repeat the raw-versus-interaction comparison at a
substantially larger width and depth, with both arms again parameter-matched
within each capacity. The quantity of interest is not either arm's absolute
score but whether the \emph{gap} between them narrows as capacity grows. A gap
that persists indicates that the operator supplies something the encoder does
not otherwise recover; a gap that closes would indicate that it merely
accelerates learning of a representation a larger model finds unaided.

\subsection{Pretraining}

\pending{The pretraining corpus and protocol are specified but not yet run. The
design prioritises corpus diversity over volume, and evaluates against a
supervised-from-scratch reference at matched budget rather than against other
pretraining variants, since our present evidence is that on tasks where
supervised training already succeeds, single-corpus masked pretraining does not
improve on no pretraining at all.}

\section{Limitations}

\pending{Limitations arising from the empirical results are deferred until the
task suite is fixed.} Two limitations are properties of the operator itself and
hold regardless of the evaluation. First, the
learned band edges are unconstrained and could in principle cross, which would
invalidate the slow-to-fast reading of $i<j$; final centre frequencies should be
inspected before any physiological claim. A one-second patch contains few cycles
of the slowest band, so the coupling estimate is noisy at that end. Finally,
statistical PAC can arise from non-sinusoidal waveform shape or sharp transients
\citep{kramer2008spurious,cole2017waveform,aru2015untangling}; we therefore call
the learned quantity a phase--amplitude dependency and make no claim that it
identifies coupling between independent oscillators.

\section{Conclusion}

A prior that can be bypassed will be. Making phase--amplitude coupling
constitutive of the token rather than a modifier of attention converts a
repeatedly failing architectural idea into a working one, and the resulting
operator is exactly invariant to the arbitrary phase reference. Parameter-matched
controls locate the gain specifically in the measured preferred phase.

\bibliography{paclock}
\bibliographystyle{iclr2026_conference}

\end{document}
